\documentclass{article}

\usepackage{
	arxiv,
	amsmath,
	amssymb,
	cleveref
}

\title{Computing the Metric Dimension of Graphs using Satisfiability Modulo Theories}

\author{
	Fikri Yudhistira Ajiwijaya \\
	Institut Teknologi Sepuluh Nopember
	\And
	Muhammad Syifa'ul Mufid \\
	Institut Teknologi Sepuluh Nopember
}

\newcommand{\combine}[7]{#1 #2 \vphantom{#6} #3 #4 #5 #6 \vphantom{#2} #7}

\begin{document}
	\maketitle

	\begin{abstract}
		...\textit{blah blah blah}...
	\end{abstract}

	\section{Introduction}

	...\textit{blah blah blah}...

	\section{Preliminaries}

	Let \(G\) be a connected graph. \\
	Let \(V = \{v_1, \cdots, v_n\}\) be the set of vertices in \(G\). \\
	Let \(E = \{e_1, \cdots, e_m\}\) be the set of edges in \(G\). \\
	Let \(W = \{w_1, \cdots, w_k\} \subseteq V\). \\
	Let \(u, v, w \in V\). \\
	Let \(e, f \in E\). \\
	Let \(e\) connects \(v\) and \(w\).

	Let \(U\) be a set with cardinality \(k\).

	Let \(C\), \(C_1\), and \(C_2\) be \(k\)-tuples. \\
	Let \(B\), \(B_1\), and \(B_2\) be Boolean k-tuples. \\
	Let \(I\), \(I_1\), and \(I_2\) be integer \(k\)-tuples.

	Let \(A\) be a matrix with \(n\) rows and \(m\) columns.

	\clearpage

	There exists a bijective function \(\mathrm{bool}\) that maps \(u \in U\) to a Boolean variable in satisfiability modulo theories. \\
	There exists a bijective function \(\mathrm{int}\) that maps \(u \in U\) to an integer variable in satisfiability modulo theories. \\
	\(X\) denotes a model given a set of constraints in satisfiability modulo theories. \\
	\(\mathcal{S}(X)\) detones the solution space of \(X\).

	\[ |U| = \text{the cardinality of } U \]
	\[ \mathcal{P}(U) = \text{the power set of } U \]
	\[ R \subseteq U,\enspace \widetilde{R} = U \backslash R \]

	\[ |C| = \text{the length of } C \]
	\[ C(i) \text{ is the element at } i \text{-th position} \]
	For \(U\) to have correspondence with \(C\) means there exists a bijective function \(\mathrm{idx} : U \to \{1, \cdots, |C|\}\).
	\[ u \in U, C(u) = C(\mathrm{idx}(u)) \]
	\[ C(\{u_1, \cdots, u_r\} \subseteq U) = (\enspace C(u_1),\enspace \cdots,\enspace C(u_r)\enspace ), \quad \text{preserving element order of } C \]

	\[ (\mathrm{True})^k \text{ is the } k \text{-tuple that contains all } \mathrm{True} \]
	\[ \neg C = (\enspace \neg C(1),\enspace \cdots,\enspace \neg C(k)\enspace ) \]
	\[ C_1 \vee C_2 = (\enspace C_1(1) \vee C_2(1),\enspace \cdots,\enspace C_1(k) \vee C_2(k)\enspace ) \]
	\[ C_1\ \ddot{=}\ C_2 = (\enspace C_1(1) = C_2(1),\enspace \cdots,\enspace C_1(k) = C_2(k)\enspace ) \]

	\[ U(B) = \{ u \in U : B(u) = \mathrm{True} \} \]
	\[ U(I) = \{ u \in U : I(u) \neq 0 \} \]
	\[ U(\{B_1, B_2, \cdots\}) = \{\enspace U(B_1),\enspace U(B_2),\enspace \cdots\enspace \} \]
	\[ U(\{I_1, I_2, \cdots\}) = \{\enspace U(I_1),\enspace U(I_2),\enspace \cdots\enspace \} \]

	\[ |A|_{\mathrm{row}} = \text{the number of rows in } A \]
	\[ |A|_{\mathrm{col}} = \text{the number of columns in } A \]
	\[ \mathcal{A} \overset{\mathrm{row}}{\in} A = \mathcal{A} \text{ is a row in } A  \text{ as an } m \text{-tuple} \]
	\[ \mathcal{A} \overset{\mathrm{col}}{\in} A = \mathcal{A} \text{ is a column in } A  \text{ as an } n \text{-tuple} \]
	\[ A(i, j) \text{ is the the element at the } i \text{-th row and } j \text{-th column} \]
	\[ A(i, \cdot) = (\enspace A(i, 1),\enspace \cdots,\enspace A(i, |A|_{\mathrm{col}})\enspace ) \]
	\[ A(\cdot, j) = (\enspace A(1, j),\enspace \cdots,\enspace A(|A|_{\mathrm{row}}, j)\enspace ) \]
	For \(S\) to have row correspondence with \(A\) means there exists a bijective function \(\mathrm{idx} : S \to \{1, \cdots, |A|_{\mathrm{row}}\}\).
	\[ s \in S, A(s, \cdot) = A(\mathrm{idx}(s), \cdot) \]
	\[
		A(\{s_1, \cdots, s_r\} \subseteq S, \cdot) = \left[
		\begin{array}{c}
			A(s_1, \cdot) \\
			\vdots \\
			A(s_r, \cdot)
		\end{array}
		\right],\quad \text{preserving row order of } A
	\]
	For \(T\) to have column correspondence with \(A\) means there exists a bijective function \(\mathrm{idx} : T \to \{1, \cdots, |A|_{\mathrm{col}}\}\).
	\[ t \in T, A(\cdot, t) = A(\cdot, \mathrm{idx}(t)) \]
	\[ A(\cdot, \{t_1, \cdots, t_r\} \subseteq T) = \left[\enspace A(\cdot, t_1)\enspace \cdots\enspace A(\cdot, t_r)\enspace \right],\quad \text{preserving column order of } A \]

	...\textit{blah blah blah}...

	\clearpage
	\section{Metric Dimension}

	\[ d(u, v) = \text{the distance from } u \text{ to } v \text{ in } G \]
	\[ r(v|W) = (\enspace d(w_1, v),\enspace d(w_2, v),\enspace \cdots,\enspace d(w_k, v)\enspace ) \]
	\[ W \text{ is a resolving set} \iff \forall u, v \in V, u \neq v \implies r(u|W) \neq r(v|W) \]
	\[ W \text{ is a non-resolving set} \iff \exists u, v \in V, u \neq v \implies r(u|W) = r(v|W) \]
	\[ W \text{ is a resolving set}, W \subseteq Z \implies Z \text{ is a resolving set} \]
	\[ W \text{ is a non-resolving set}, Z \subseteq W \implies Z \text{ is a non-resolving set} \]
	\[ \mathcal{P}(V) \text{ is partitioned into the set of resolving sets and the set of non-resolving sets} \]
	\[ \text{The metric dimension of } G \text{ is the minimum cardinality of its resolving set} \]

	\[ r(V|W) = \left[\enspace r(v_1|W)\enspace r(v_2|W)\enspace \cdots\enspace r(v_n|W)\enspace \right] \]
	\(V\) has correspondence with \(r(v|W)\). \\
	\(V\) has column correspondence \(r(V|W)\).
	\[ r(v|W) = r(v|V)(W) \]
	\[ r(v|W) = r(V|W)(\cdot, v) \]
	\(V\) has row correspondence \(r(V|V)\).
	\[ r(V|W) = r(V|V)(W, \cdot) \]
	\[ r(v|W) = r(V|V)(W, \cdot)(\cdot, v) \]
	\[ W \text{ is a resolving set} \iff \forall u, v \in V, u \neq v \implies r(V|W)(\cdot, u) \neq r(V|W)(\cdot, v) \]
	\[ W \text{ is a non-resolving set} \iff \exists u, v \in V, u \neq v \implies r(V|W)(\cdot, u) = r(V|W)(\cdot, v) \]

	\[ d(u, e) = \textrm{min}(\enspace d(u, v),\enspace d(u, w)\enspace ) \]
	\[ r(e|W) = (\enspace d(w_1, e),\enspace d(w_2, e),\enspace \cdots,\enspace d(w_k, e)\enspace ) \]
	\[ W \text{ is an edge resolving set} \iff \forall e, f \in E, e \neq f \implies r(e|W) \neq r(f|W) \]
	\[ W \text{ is an edge non-resolving set} \iff \exists e, f \in E, e \neq f \implies r(e|W) = r(f|W) \]
	\[ W \text{ is an edge resolving set}, W \subseteq Z \implies Z \text{ is an edge resolving set} \]
	\[ W \text{ is an edge non-resolving set}, Z \subseteq W \implies Z \text{ is an edge non-resolving set} \]
	\[ \mathcal{P}(V) \text{ is partitioned into the set of edge resolving sets and the set of edge non-resolving sets} \]
	\[ \text{The edge metric dimension of } G \text{ is the minimum cardinality of its edge resolving set} \]

	\[ r(E|W) = \left[\enspace r(e_1|W)\enspace r(e_2|W)\enspace \cdots\enspace r(e_m|W)\enspace \right] \]
	\(V\) has correspondence with \(r(v|W)\). \\
	\(E\) has column correspondence \(r(E|W)\).
	\[ r(e|W) = r(e|V)(W) \]
	\[ r(e|W) = r(E|W)(\cdot, e) \]
	\(V\) has row correspondence \(r(E|V)\).
	\[ r(E|W) = r(E|V)(W, \cdot) \]
	\[ r(e|W) = r(E|V)(W, \cdot)(\cdot, e) \]
	\[ W \text{ is an edge resolving set} \iff \forall e, f \in E, e \neq f \implies r(E|W)(\cdot, e) \neq r(E|W)(\cdot, f) \]
	\[ W \text{ is an edge non-resolving set} \iff \exists e, f \in E, e \neq f \implies r(E|W)(\cdot, e) = r(E|W)(\cdot, f) \]

	...\textit{blah blah blah}...

	\clearpage
	\section{Distance Matrix}

	\[
		D^{(V)} = \left[
		\begin{array}{cccc}
			d(v_1, v_1) & d(v_1, v_2) & \cdots & d(v_1, v_n) \\
			d(v_2, v_1) & d(v_2, v_2) & \cdots & d(v_2, v_n) \\
			\vdots & \vdots & \ddots & \vdots \\
			d(v_n, v_1) & d(v_n, v_2) & \cdots & d(v_n, v_n) \\
		\end{array}
		\right]
	\]
	\[ D^{(V)} = \left[\enspace r(v_1|V)\enspace r(v_2|V)\enspace \cdots\enspace r(v_n|V)\enspace \right] \]
	\[ D^{(V)} = r(V|V) \]
	\(D^{(V)}\) has a dimension of \(|V| \times |V|\). \\
	\(V\) has row correspondence with \(D^{(V)}\). \\
	\(V\) has column correspondence with \(D^{(V)}\).
	\[ r(v|V) = D^{(V)}(\cdot, v) \]
	\[ r(V|W) = D^{(V)}(W, \cdot) \]
	\[ r(v|W) = D^{(V)}(\cdot, v)(W) \]
	\[ r(v|W) = D^{(V)}(W, \cdot)(\cdot, v) \]
	\[ W \text{ is a resolving set} \iff \forall u, v \in V, u \neq v \implies D^{(V)}(W, \cdot)(\cdot, u) \neq D^{(V)}(W, \cdot)(\cdot, v) \]
	\[ W \text{ is a non-resolving set} \iff \exists u, v \in V, u \neq v \implies D^{(V)}(W, \cdot)(\cdot, u) = D^{(V)}(W, \cdot)(\cdot, v) \]

	\[
		D^{(E)} = \left[
		\begin{array}{cccc}
			d(v_1, e_1) & d(v_1, e_2) & \cdots & d(v_1, e_m) \\
			d(v_2, e_1) & d(v_2, e_2) & \cdots & d(v_2, e_m) \\
			\vdots & \vdots & \ddots & \vdots \\
			d(v_n, e_1) & d(v_n, e_2) & \cdots & d(v_n, e_m) \\
		\end{array}
		\right]
	\]
	\[ D^{(E)} = \left[\enspace r(e_1|V)\enspace r(e_2|V)\enspace \cdots\enspace r(e_m|V)\enspace \right] \]
	\[ D^{(E)} = r(E|V) \]
	\(D^{(E)}\) has a dimension of \(|V| \times |E|\). \\
	\(V\) has row correspondence with \(D^{(E)}\). \\
	\(E\) has column correspondence with \(D^{(E)}\).
	\[ r(e|V) = D^{(E)}(\cdot, e) \]
	\[ r(E|W) = D^{(E)}(W, \cdot) \]
	\[ r(e|W) = D^{(E)}(\cdot, e)(W) \]
	\[ r(e|W) = D^{(E)}(W, \cdot)(\cdot, e) \]
	\[ W \text{ is an edge resolving set} \iff \forall e, f \in E, e \neq f \implies D^{(E)}(W, \cdot)(\cdot, e) \neq D^{(E)}(W, \cdot)(\cdot, f) \]
	\[ W \text{ is an edge non-resolving set} \iff \exists e, f \in E, e \neq f \implies D^{(E)}(W, \cdot)(\cdot, e) = D^{(E)}(W, \cdot)(\cdot, f) \]

	...\textit{blah blah blah}...

	\clearpage
	\section{Distance Similarity Matrix}

	\(V\) has correspondence with \(r(u|V)\ \ddot{=}\ r(v|V)\).
	\[ r(u|W)\ \ddot{=}\ r(v|W) = (r(u|V)\ \ddot{=}\ r(v|V))(W) \]

	\[ B^{(V)} = \left[\enspace r(v_i|V)\ \ddot{=}\ r(v_j|V)\enspace \cdots\enspace \right],\quad 1 \leq i < j \leq |V| \]
	\[ B^{(V)} = \left[\enspace D^{(V)}(\cdot, v_i)\ \ddot{=}\ D^{(V)}(\cdot, v_j) \enspace \cdots\enspace \right],\quad 1 \leq i < j \leq |V| \]
	\(V\) has row correspondence with \(B^{(V)}\).
	\[ \forall u, v \in V, u \neq v,\enspace \exists C \overset{\text{col}}{\in} B^{(V)},\enspace C = r(u|V)\ \ddot{=}\ r(v|V) \]
	\[ W \text{ is a resolving set} \iff \forall C \overset{\text{col}}{\in} B^{(V)},\enspace C(W) \neq (\mathrm{True})^{|W|} \]
	\[ W \text{ is a non-resolving set} \iff \exists C \overset{\text{col}}{\in} B^{(V)},\enspace C(W) = (\mathrm{True})^{|W|} \]

	\vspace{2em}

	\(V\) has correspondence with \(r(e|V)\ \ddot{=}\ r(f|V)\).
	\[ r(e|W)\ \ddot{=}\ r(f|W) = ((e|V)\ \ddot{=}\ r(f|V))(W) \]

	\[ B^{(E)} = \left[\enspace r(e_i|V)\ \ddot{=}\ r(e_j|V)\enspace \cdots\enspace \right],\quad 1 \leq i < j \leq |E| \]
	\[ B^{(E)} = \left[\enspace D^{(E)}(\cdot, e_i)\ \ddot{=}\ D^{(E)}(\cdot, e_j) \enspace \cdots\enspace \right],\quad 1 \leq i < j \leq |E| \]
	\(V\) has row correspondence with \(B^{(E)}\).
	\[ \forall e, f \in E, e \neq f,\enspace \exists C \overset{\mathrm{col}}{\in} B^{(E)},\enspace C = r(e|V)\ \ddot{=}\ r(f|V) \]
	\[ W \text{ is an edge resolving set} \iff \forall C \overset{\mathrm{col}}{\in} B^{(E)},\enspace C(W) \neq (\mathtt{True})^{|W|} \]
	\[ W \text{ is an edge non-resolving set} \iff \exists C \overset{\mathrm{col}}{\in} B^{(E)},\enspace C(W) = (\mathtt{True})^{|W|} \]

	...\textit{blah blah blah}...

	\clearpage
	\section{Reduced Distance Similarity Matrix}

	\[ C \overset{\mathrm{col}}{\in} B^{(V)},\enspace V(C) \text{ is a non-resolving set} \]
	\[ C, \mathcal{C} \overset{\mathrm{col}}{\in} B^{(V)},\enspace V(C) \subseteq V(\mathcal{C}) \iff C \vee \mathcal{C} = \mathcal{C} \]
	\(P^{(V)}\) is the submatrix of \(B^{(V)}\) obtained by eliminating all strictly dominated columns and reducing multiple identical columns to a single instance.
	\[ W \text{ is a resolving set} \iff \forall C \overset{\mathrm{col}}{\in} P^{(V)},\enspace C(W) \neq (\mathrm{True})^{|W|} \]
	\[ W \text{ is a non-resolving set} \iff \exists C \overset{\mathrm{col}}{\in} P^{(V)},\enspace C(W) = (\mathrm{True})^{|W|} \]
	\[ C \overset{\mathrm{col}}{\in} P^{(V)},\enspace V(C) \text{ is a non-resolving set} \]

	\[ C \overset{\mathrm{col}}{\in} B^{(E)},\enspace V(C) \text{ is an edge non-resolving set} \]
	\[ C, \mathcal{C} \overset{\mathrm{col}}{\in} B^{(E)},\enspace V(C) \subseteq V(\mathcal{C}) \iff C \vee \mathcal{C} = \mathcal{C} \]
	\(P^{(E)}\) is the submatrix of \(B^{(E)}\) obtained by eliminating all strictly dominated columns and reducing multiple identical columns to a single instance.
	\[ W \text{ is an edge resolving set} \iff \forall C \overset{\mathrm{col}}{\in} P^{(E)},\enspace C(W) \neq (\mathtt{True})^{|W|} \]
	\[ W \text{ is an edge non-resolving set} \iff \exists C \overset{\mathrm{col}}{\in} P^{(E)},\enspace C(W) = (\mathtt{True})^{|W|} \]
	\[ C \overset{\mathrm{col}}{\in} P^{(E)},\enspace V(C) \text{ is an edge non-resolving set} \]

	...\textit{blah blah blah}...

	\clearpage
	\section{Satisfiability Modulo Theories}

	\[ \mathcal{S}(X) = \emptyset \iff \text{unsatisfiable} \]
	\[ \mathcal{S}(X) \neq \emptyset \iff \text{satisfiable} \]

	\[ \text{Given zero constraints},\enspace V(\mathcal{S}(X)) = \mathcal{P}(V) \]

	\[ W \subseteq V,\qquad \combine{\left(}{\bigwedge \limits_{\displaystyle v \in W} \mathrm{bool}(v)}{\right)}{\bigwedge}{\left(}{\bigwedge \limits_{\displaystyle v \in \widetilde{W}} \neg \mathrm{bool}(v)}{\right)}\enspace \implies\enspace V(X) = W \]

	\[ W \subseteq V,\quad \bigwedge \limits_{\displaystyle v \in \widetilde{W}} \neg \mathrm{bool}(v) \implies V(\mathcal{S}(X)) = \mathcal{P}(W) \]

	...\textit{blah blah blah}...

	\clearpage
	\section{Boolean Satisfiability}

	\[ W \text{ is a non-resolving set},\quad \bigwedge \limits_{\displaystyle v \in \widetilde{W}} \neg \mathrm{bool}(v) \implies V(X) \text{ is a non-resolving set} \]

	\[ C \overset{\mathrm{col}}{\in} P^{(V)},\quad \bigwedge \limits_{\displaystyle v \in V(\neg C)} \neg \mathrm{bool}(v) \implies V(X) \text{ is a non-resolving set} \]

	\[ \bigvee \limits_{\displaystyle C \overset{\mathrm{col}}{\in} P^{(V)}} \left( \bigwedge \limits_{\displaystyle v \in V(\neg C)} \neg \mathrm{bool}(v) \right)\enspace \implies\enspace V(X) \text{ is a non-resolving set} \]

	\[ \bigvee \limits_{\displaystyle C \overset{\mathrm{col}}{\in} P^{(V)}} \left( \bigwedge \limits_{\displaystyle v \in V(\neg C)} \neg \mathrm{bool}(v) \right)\enspace \implies\enspace V(\mathcal{S}(X)) \text{ is the non-resolving set partition} \]

	\[ \bigwedge \limits_{\displaystyle C \overset{\mathrm{col}}{\in} P^{(V)}} \left( \bigvee \limits_{\displaystyle v \in V(\neg C)} \mathrm{bool}(v) \right)\enspace \implies\enspace V(\mathcal{S}(X)) \text{ is the resolving set partition} \]

	\[ \bigwedge \limits_{\displaystyle C \overset{\mathrm{col}}{\in} P^{(V)}} \left( \bigvee \limits_{\displaystyle v \in V(\neg C)} \mathrm{bool}(v) \right)\enspace \implies\enspace V(X) \text{ is a resolving set} \]

	\[ W \text{ is an edge non-resolving set},\quad \bigwedge \limits_{\displaystyle v \in \widetilde{W}} \neg \mathrm{bool}(v) \implies V(X) \text{ is an edge non-resolving set} \]

	\[ C \overset{\mathrm{col}}{\in} P^{(E)},\quad \bigwedge \limits_{\displaystyle v \in V(\neg C)} \neg \mathrm{bool}(v) \implies V(X) \text{ is an edge non-resolving set} \]

	\[ \bigvee \limits_{\displaystyle C \overset{\mathrm{col}}{\in} P^{(E)}} \left( \bigwedge \limits_{\displaystyle v \in V(\neg C)} \neg \mathrm{bool}(v) \right)\enspace \implies\enspace V(X) \text{ is an edge non-resolving set} \]

	\[ \bigvee \limits_{\displaystyle C \overset{\mathrm{col}}{\in} P^{(E)}} \left( \bigwedge \limits_{\displaystyle v \in V(\neg C)} \neg \mathrm{bool}(v) \right)\enspace \implies\enspace V(\mathcal{S}(X)) \text{ is the edge non-resolving set partition} \]

	\[ \bigwedge \limits_{\displaystyle C \overset{\mathrm{col}}{\in} P^{(E)}} \left( \bigvee \limits_{\displaystyle v \in V(\neg C)} \mathrm{bool}(v) \right)\enspace \implies\enspace V(\mathcal{S}(X)) \text{ is the edge resolving set partition} \]

	\[ \bigwedge \limits_{\displaystyle C \overset{\mathrm{col}}{\in} P^{(E)}} \left( \bigvee \limits_{\displaystyle v \in V(\neg C)} \mathrm{bool}(v) \right)\enspace \implies\enspace V(X) \text{ is an edge resolving set} \]

	...\textit{blah blah blah}...

	\clearpage
	\subsection{Transformation}

	\[
	\begin{array}{l}
		\displaystyle \bigwedge \limits_{\displaystyle C \overset{\mathrm{col}}{\in} P} \left( \combine{\left(}{\mathrm{False}}{\right)}{\bigvee}{\left(}{\bigvee \limits_{\displaystyle v \in V(\neg C)} \mathrm{bool}(v)}{\right)} \right) \\ [24pt]

		\displaystyle\qquad = \bigwedge \limits_{\displaystyle C \overset{\mathrm{col}}{\in} P} \left( \combine{\left(}{\combine{\left(}{\bigvee \limits_{\displaystyle v \in V(C)} \mathrm{bool}(v)}{\right)}{\bigwedge}{\left(}{\neg \bigvee \limits_{\displaystyle v \in V(C)} \mathrm{bool}(v)}{\right)}}{\right)}{\bigvee}{\left(}{\bigvee \limits_{\displaystyle v \in V(\neg C)} \mathrm{bool}(v)}{\right)} \right) \\ [24pt]

		\displaystyle\qquad = \bigwedge \limits_{\displaystyle C \overset{\mathrm{col}}{\in} P} \left( \begin{array}{c}
			\left( \combine{\left(}{\bigvee \limits_{\displaystyle v \in V(C)} \mathrm{bool}(v)}{\right)}{\bigvee}{\left(}{\bigvee \limits_{\displaystyle v \in V(\neg C)} \mathrm{bool}(v)}{\right)} \right) \\ [24pt]
			\bigwedge \\ [16pt]
			\left( \combine{\left(}{\neg \bigvee \limits_{\displaystyle v \in V(C)} \mathrm{bool}(v)}{\right)}{\bigvee}{\left(}{\bigvee \limits_{\displaystyle v \in V(\neg C)} \mathrm{bool}(v)}{\right)} \right)
		\end{array}
		\right) \\ [58pt]

		\displaystyle\qquad = \bigwedge \limits_{\displaystyle C \overset{\mathrm{col}}{\in} P} \left( \combine{\left(}{\bigvee \limits_{\displaystyle v \in V} \mathrm{bool}(v)}{\right)}{\bigwedge}{\left(}{\bigvee \limits_{\displaystyle v \in V(C)} \mathrm{bool}(v)\ \implies \bigvee \limits_{\displaystyle v \in V(\neg C)} \mathrm{bool}(v)}{\right)} \right) \\ [24pt]

		\displaystyle\qquad = \combine{\left(}{\bigvee \limits_{\displaystyle v \in V} \mathrm{bool}(v)}{\right)}{\bigwedge}{\left(}{\bigwedge \limits_{\displaystyle C \overset{\mathrm{col}}{\in} P} \left( \bigvee \limits_{\displaystyle v \in V(C)} \mathrm{bool}(v)\ \implies \bigvee \limits_{\displaystyle v \in V(\neg C)} \mathrm{bool}(v) \right)}{\right)}
	\end{array}
	\]

	...\textit{blah blah blah}...

	\[
	\begin{array}{c}
		\combine{\left(}{\bigvee \limits_{\displaystyle v \in V} \mathrm{bool}(v)}{\right)}{\bigwedge}{\left(}{\bigwedge \limits_{\displaystyle C \overset{\mathrm{col}}{\in} P^{(V)}} \left( \bigvee \limits_{\displaystyle v \in V(C)} \mathrm{bool}(v)\ \implies \bigvee \limits_{\displaystyle v \in V(\neg C)} \mathrm{bool}(v) \right)}{\right)} \\ [24pt]
		\implies\enspace V(X) \text{ is a resolving set}
	\end{array}
	\]

	\[
	\begin{array}{c}
		\combine{\left(}{\bigvee \limits_{\displaystyle v \in V} \mathrm{bool}(v)}{\right)}{\bigwedge}{\left(}{\bigwedge \limits_{\displaystyle C \overset{\mathrm{col}}{\in} P^{(E)}} \left( \bigvee \limits_{\displaystyle v \in V(C)} \mathrm{bool}(v)\ \implies \bigvee \limits_{\displaystyle v \in V(\neg C)} \mathrm{bool}(v) \right)}{\right)} \\ [24pt]
		\implies\enspace V(X) \text{ is an edge resolving set}
	\end{array}
	\]

	...\textit{blah blah blah}...

	\clearpage
	\subsection{Cardinality}

	\[ \cdots \bigwedge \mathrm{AtLeast}(k) \bigwedge \mathrm{AtMost}(k) \implies V(\mathcal{S}(X)) \text{ is the set of resolving sets of cardinality } k \]

	...\textit{blah blah blah}...

	It is possible to find the true minimum iteratively in descending order from the top, (\(k, k - 1, k - 2, \cdots\)), due to monotonicity.

	...\textit{blah blah blah}...

	\clearpage
	\section{Linear Integer Arithmetic}

	\[ \combine{\left(}{\bigwedge \limits_{\displaystyle v \in V} 0 \leq \mathrm{int}(v) \leq 1}{\right)}{\bigwedge}{\left(}{\bigwedge \limits_{\displaystyle C \overset{\mathrm{col}}{\in} P^{(V)}} \left( \sum \limits_{\displaystyle v \in V(\neg C)} \mathrm{int}(v)\enspace \geq\enspace 1 \right)}{\right)}\enspace \implies\enspace V(X) \text{ is a resolving set} \]

	\[ \combine{\left(}{\bigwedge \limits_{\displaystyle v \in V} 0 \leq \mathrm{int}(v) \leq 1}{\right)}{\bigwedge}{\left(}{\bigwedge \limits_{\displaystyle C \overset{\mathrm{col}}{\in} P^{(E)}} \left( \sum \limits_{\displaystyle v \in V(\neg C)} \mathrm{int}(v)\enspace \geq\enspace 1 \right)}{\right)}\enspace \implies\enspace V(X) \text{ is an edge resolving set} \]

	...\textit{blah blah blah}...

	\[ \cdots \bigwedge \sum \limits_{\displaystyle v \in V} \mathrm{int}(v) = k \implies V(\mathcal{S}(X)) \text{ is the set of resolving sets of cardinality } k \]

	...\textit{blah blah blah}...

	\[ \cdots \bigwedge \left( \bigwedge \limits_{\displaystyle C \overset{\mathrm{col}}{\in} P^{(V)}} \left( \sum \limits_{\displaystyle v \in V(\neg C)} \mathrm{int}(v)\enspace \leq\enspace k - 1 \right) \right)\enspace \implies\enspace V(X) \text{ is a resolving set} \]

	\[ \cdots \bigwedge \left( \bigwedge \limits_{\displaystyle C \overset{\mathrm{col}}{\in} P^{(E)}} \left( \sum \limits_{\displaystyle v \in V(\neg C)} \mathrm{int}(v)\enspace \leq\enspace k - 1 \right) \right)\enspace \implies\enspace V(X) \text{ is an edge resolving set} \]

	\clearpage
	\section{Pseudo Boolean}

	\[ \bigwedge \limits_{\displaystyle C \overset{\mathrm{col}}{\in} P^{(V)}} \left( \sum \limits_{\displaystyle v \in V(\neg C)} \mathrm{bool}(v)\enspace \leq\enspace k - 1 \right)\enspace \implies\enspace V(X) \text{ is a resolving set} \]

	\[ \bigwedge \limits_{\displaystyle C \overset{\mathrm{col}}{\in} P^{(E)}} \left( \sum \limits_{\displaystyle v \in V(\neg C)} \mathrm{bool}(v)\enspace \leq\enspace k - 1 \right)\enspace \implies\enspace V(X) \text{ is an edge resolving set} \]

	...\textit{blah blah blah}...

	\clearpage
	\section{Experiments and Results}

	...\textit{blah blah blah}...

	\section{Conclusions and Future Works}

	...\textit{blah blah blah}...
\end{document}
